% GENERATED FILE. Edit canonical lessons, then run npm run generate:books.
% canonical-source-hash: fnv1a64:dec8452c1eb2d9f3
% canonical-lessons: HI-C197-kaksa, HI-C197-pariksa, HI-C197-path, HI-C197-abhyas, HI-C197-ank

\chapter{A class, An examination, A lesson, Practice, A mark}
\label{ch:hi-a-class-an-examination-a-lesson-practice-a-mark}
\hlchaptermodality{197}

\begin{chapteropening}
\textbf{By the end of this chapter:} \emph{I can say a class, an examination, a lesson, practice and a mark.}

Complete the last lesson of chapter 197: I can say a class, an examination, a lesson, practice and a mark.
\end{chapteropening}

\section[\texorpdfstring{\hi{कक्षा} (kakṣā)}{कक्षा (kakṣā)}]{\hi{कक्षा} (kakṣā) — a class, a classroom}

\label{lesson:HI-C197-kaksa}



\begin{quote}
Before the new one: say the Hindi for a servant, then the Hindi for a thief.
\end{quote}

\subsection*{You'll want to know: \hi{कक्षा}}
\textbf{\hi{कक्षा}} — \emph{kakṣā} — \textquotedblleft{}a class, a classroom\textquotedblright{}.

\begin{cousinweb}
\textbf{\hi{कक्षा}} (\emph{kakṣā}) is Sanskrit \emph{kakṣā}, \textquotedblleft{}girdle\textquotedblright{}, taken over whole as a learned loan, a tatsama (CDIAL 2592). Hindi uses it for a classroom, and so for a class; CDIAL does not trace that sense.
\end{cousinweb}

\begin{grammarlens}[title={What you've built}]
One more word for this chapter.
\end{grammarlens}

\subsection*{Your turn}
\begin{itemize}
\raggedright
  \item \emph{Say it:} \emph{kakṣā}
  \item \emph{Say it:} \emph{kakṣā}, once more
  \item \emph{Say it:} say \emph{kakṣā}
  \item \emph{Recall:} say the Hindi for a servant, then the Hindi for a thief, then say \emph{kakṣā} again
\end{itemize}

\begin{tcolorbox}[breakable,colback=teal!4,colframe=teal!35!black,title={Before you move on}]
What does \hi{कक्षा} mean? (\textquotedblleft{}A class, a classroom\textquotedblright{}.) Say it once more.
\end{tcolorbox}
\section[\texorpdfstring{\hi{परीक्षा} (parīkṣā)}{परीक्षा (parīkṣā)}]{\hi{परीक्षा} (parīkṣā) — an examination}

\label{lesson:HI-C197-pariksa}



\begin{quote}
Before the new one: say the Hindi for a thief, then the Hindi for a class.
\end{quote}

\subsection*{You'll want to know: \hi{परीक्षा}}
\textbf{\hi{परीक्षा}} — \emph{parīkṣā} — \textquotedblleft{}an examination\textquotedblright{}.

\begin{cousinweb}
\textbf{\hi{परीक्षा}} (\emph{parīkṣā}) is Sanskrit \emph{parīkṣā}, \textquotedblleft{}inspection, test\textquotedblright{}, taken over whole as a learned loan, a tatsama (CDIAL 7904). CDIAL lists the inherited form of the same word as \emph{parakh}, \textquotedblleft{}examination\textquotedblright{}.
\end{cousinweb}

\begin{grammarlens}[title={What you've built}]
One more word for this chapter.
\end{grammarlens}

\subsection*{Your turn}
\begin{itemize}
\raggedright
  \item \emph{Say it:} \emph{parīkṣā}
  \item \emph{Say it:} \emph{parīkṣā}, once more
  \item \emph{Say it:} say \emph{parīkṣā}
  \item \emph{Recall:} say the Hindi for a thief, then the Hindi for a class, then say \emph{parīkṣā} again
\end{itemize}

\begin{tcolorbox}[breakable,colback=teal!4,colframe=teal!35!black,title={Before you move on}]
What does \hi{परीक्षा} mean? (\textquotedblleft{}An examination\textquotedblright{}.) Say it once more.
\end{tcolorbox}
\section[\texorpdfstring{\hi{पाठ} (pāṭh)}{पाठ (pāṭh)}]{\hi{पाठ} (pāṭh) — a lesson}

\label{lesson:HI-C197-path}



\begin{quote}
Before the new one: say the Hindi for a class, then the Hindi for an examination.
\end{quote}

\subsection*{You'll want to know: \hi{पाठ}}
\textbf{\hi{पाठ}} — \emph{pāṭh} — \textquotedblleft{}a lesson\textquotedblright{}.

\begin{cousinweb}
\textbf{\hi{पाठ}} (\emph{pāṭh}) is Sanskrit \emph{pāṭha}, \textquotedblleft{}text of a book; study\textquotedblright{}, taken over whole as a learned loan, a tatsama (CDIAL 8035).
\end{cousinweb}

\begin{grammarlens}[title={What you've built}]
One more word for this chapter.
\end{grammarlens}

\subsection*{Your turn}
\begin{itemize}
\raggedright
  \item \emph{Say it:} \emph{pāṭh}
  \item \emph{Say it:} \emph{pāṭh}, once more
  \item \emph{Say it:} say \emph{pāṭh}
  \item \emph{Recall:} say the Hindi for a class, then the Hindi for an examination, then say \emph{pāṭh} again
\end{itemize}

\begin{tcolorbox}[breakable,colback=teal!4,colframe=teal!35!black,title={Before you move on}]
What does \hi{पाठ} mean? (\textquotedblleft{}A lesson\textquotedblright{}.) Say it once more.
\end{tcolorbox}
\section[\texorpdfstring{\hi{अभ्यास} (abhyās)}{अभ्यास (abhyās)}]{\hi{अभ्यास} (abhyās) — practice, an exercise}

\label{lesson:HI-C197-abhyas}



\begin{quote}
Before the new one: say the Hindi for an examination, then the Hindi for a lesson.
\end{quote}

\subsection*{You'll want to know: \hi{अभ्यास}}
\textbf{\hi{अभ्यास}} — \emph{abhyās} — \textquotedblleft{}practice, an exercise\textquotedblright{}.

\begin{cousinweb}
\textbf{\hi{अभ्यास}} (\emph{abhyās}) is Sanskrit \emph{abhyāsa}, \textquotedblleft{}repetition, practice\textquotedblright{}, taken over whole as a learned loan, a tatsama.
\end{cousinweb}

\begin{grammarlens}[title={What you've built}]
One more word for this chapter.
\end{grammarlens}

\subsection*{Your turn}
\begin{itemize}
\raggedright
  \item \emph{Say it:} \emph{abhyās}
  \item \emph{Say it:} \emph{abhyās}, once more
  \item \emph{Say it:} say \emph{abhyās}
  \item \emph{Recall:} say the Hindi for an examination, then the Hindi for a lesson, then say \emph{abhyās} again
\end{itemize}

\begin{tcolorbox}[breakable,colback=teal!4,colframe=teal!35!black,title={Before you move on}]
What does \hi{अभ्यास} mean? (\textquotedblleft{}Practice, an exercise\textquotedblright{}.) Say it once more.
\end{tcolorbox}
\section[\texorpdfstring{\hi{अंक} (aṅk)}{अंक (aṅk)}]{\hi{अंक} (aṅk) — a mark, a number}

\label{lesson:HI-C197-ank}



\begin{quote}
Before the new one: say the Hindi for a lesson, then the Hindi for practice.
\end{quote}

\subsection*{You'll want to know: \hi{अंक}}
\textbf{\hi{अंक}} — \emph{aṅk} — \textquotedblleft{}a mark, a number\textquotedblright{}.

\begin{cousinweb}
\textbf{\hi{अंक}} (\emph{aṅk}) is Sanskrit \emph{aṅka}, \textquotedblleft{}hook; mark, sign\textquotedblright{}, taken over whole as a learned loan, a tatsama (CDIAL 100). CDIAL lists the inherited form as \emph{āṁk}, \textquotedblleft{}mark, spot\textquotedblright{}.
\end{cousinweb}

\begin{grammarlens}[title={What you've built}]
One more word for this chapter.
\end{grammarlens}

\subsection*{Your turn}
\begin{itemize}
\raggedright
  \item \emph{Say it:} \emph{aṅk}
  \item \emph{Say it:} \emph{aṅk}, once more
  \item \emph{Say it:} say \emph{aṅk}
  \item \emph{Recall:} say the Hindi for a lesson, then the Hindi for practice, then say \emph{aṅk} again
\end{itemize}

\begin{tcolorbox}[breakable,colback=teal!4,colframe=teal!35!black,title={Before you move on}]
What does \hi{अंक} mean? (\textquotedblleft{}A mark, a number\textquotedblright{}.) Say it once more.
\end{tcolorbox}
